\documentclass{article}
\usepackage[T1]{fontenc}
\usepackage{lmodern}
\usepackage{graphicx}
\usepackage{amsmath,amssymb}
\usepackage[a4paper,margin=2cm]{geometry}
\usepackage[absolute,overlay]{textpos}
\setlength{\TPHorizModule}{1mm}
\setlength{\TPVertModule}{1mm}
\usepackage[indonesian]{babel}
\usepackage{hyperref}
\setlength{\emergencystretch}{2em}
% unit: Halaman 26

\begin{document}

% === Halaman 26 (python/native) ===

• Contoh: Tentukan semua titik P(x,y) pada kurva eliptik 


y2  x3 + x + 6  (mod 11)  

    dengan x dan y didefinisikan di dalam GF(11)

 Jawab:

 x = 0 → y2  6  (mod 11) → tidak ada nilai y yang memenuhi

 x = 1 → y2   8  (mod 11) → tidak ada nilai y yang memenuhi

 x = 2 → y2  16  (mod 11)  5 (mod 11) → y1 = 4 dan y2 = 7






         →titik P(2,4)  dan P’(2, 7)

 x = 3 → y2  36  (mod 11)  3  (mod 11) → y1 = 5 dan y2 = 6






         →titik P(3,5)  dan P’(3, 6)

26

\end{document}
